\documentclass[11pt]{letter}

\usepackage[margin=0.75in]{geometry}
\usepackage{setspace}
\usepackage{hyperref}
\usepackage{xcolor}
\usepackage{enumitem}
\usepackage{parskip}

\hypersetup{
    colorlinks=true,
    urlcolor=black,
    linkcolor=black
}

\setlength{\parindent}{0pt}
\setstretch{1.05}
\date{}
\begin{document}

\begin{letter}{}

\opening{Dear Hiring Committee,}

I saw the job post on \textbf{Lenovo's careers website} and thought I'd be a good fit for it.

I am a Master's student in Computer Science at the University of Cincinnati, graduating December 2026, looking for a role as an \textbf{AI Enterprise Engineer} at \textbf{Lenovo}. With hands on experience building backend services, REST APIs, cloud platforms, and reliable deployment systems, I am excited about the opportunity to contribute to \textbf{Lenovo} and support the \textbf{Qira enterprise platform team} in building enterprise capabilities spanning identity, administration, governance, and knowledge retrieval.

\textbf{Here are reasons that make me a great fit for the role:}

\begin{itemize}[leftmargin=*, itemsep=0.45em]
    \item \textbf{Professional software engineering experience with high-level languages, APIs, and databases:} At Siemens Industry Software, I built backend automation services in Python (Boto3) integrating with AWS APIs, replacing legacy Ruby scripts and cutting deployment time by 40\% across enterprise environments with PostgreSQL and MySQL backed systems.

    \item \textbf{Cloud computing and enterprise platform experience:} I migrated 40+ production environments to Kubernetes using ArgoCD, Helm, GitLab CI/CD, and HashiCorp Vault, strengthening identity and secrets management, and executed 44 production upgrades in 4 months across AWS, Linux, and enterprise database services.

    \item \textbf{Collaborative software development and operational excellence:} I designed an Ansible based Blue-Green deployment framework that reduced production upgrade downtime by 50\%, and served as Quality Coach and SonarQube Specialist, increasing automated code coverage from 35\% to 78\% while mentoring 5 engineers.

    \item \textbf{Exposure to AI applications and LLMs:} In my graduate research, I designed a memory-augmented LLM pipeline using Google Gemini for Android bug reproduction, achieving a 25.7x latency reduction over a published baseline. I also built a 17-tool MCP server exposing structured LLM tool-calling for enterprise data workflows in my PolicyPilot project.
\end{itemize}

I am available for full-time work after December 2026. As a Master's student at the University of Cincinnati and someone who values professionalism and growth, I am enthusiastic about contributing to Lenovo's mission of delivering Smarter Technology for All.

Please find my CV attached below. I look forward to hearing from you. I would be excited to help build the enterprise AI capabilities that let Qira's customers deploy AI securely and at scale.

{\setlength{\parskip}{0pt}
Regards,\\
Tanmay Bhuskute,\\
bhuskutv@mail.uc.edu $\mid$ tanmay.v.bhuskute@gmail.com $\mid$ +1 7328583027,\\
3424 Brookline Ave, Flat 19, Cincinnati, OH 45220.
}

\end{letter}

\end{document}
